% Set configurable latex parameters for this paper.
\input{configure_coulomb.tex}
\shortauthors{Irwin}
\shorttitle{The Coulomb Treatment}
\begin{document}

\title{The FreeEOS Code for Calculating the Equation of State for Stellar
Interiors III: The Treatment of the Coulomb Correction}

\author{Alan W. Irwin}
\affil{Department of Physics and Astronomy, University of Victoria,\\
P.O. Box 3055, Victoria, British Columbia, Canada, V8W 3P6\\
Electronic mail: irwin@beluga.phys.uvic.ca}

\begin{abstract}
This paper describes how the Coulomb correction
is implemented in FreeEOS (\url{http://freeeos.sourceforge.net/}), a
software package for rapidly calculating the equation of state for physical
conditions in stellar interiors.
\end{abstract}
\keywords{Coulomb interaction --- equation of state --- stellar interiors
--- stellar evolution}

\section{Introduction}

FreeEOS (\url{http://freeeos.sourceforge.net/}) is a software package for
rapidly calculating the equation of state (hereafter, EOS) for stellar
conditions, and a series of papers is being prepared that describe its
implementation. Paper~I (Irwin 2004a) describes the Fermi-Dirac integral
approximations. Paper~II (Irwin 2004b) describes the efficient method of
solution that delivers thermodynamically consistent results of high
numerical quality that are in excellent agreement with OPAL (Rogers \&
Nayfonov 2002) results for the solar case. Here in this third paper in the
series, I present the implementation of the Coulomb correction. This
non-ideal correction accounts for the tendency of oppositely charged,
unbound particles to correlate with each other because of the attractive
Coulomb forces between them. Because of this net attraction between the
correlated particles, less pressure is generally required to confine the gas
to a particular volume for a given temperature. An accurate treatment of the
Coulomb correction is essential since this correction is one of the most
important non-ideal corrections to the EOS calculation for stellar-interior
conditions.

The remainder of this paper is organized as follows:
Section~\ref{coulomb_sect} presents our recommended model for the Coulomb
free energy as well as various other models I have programmed as a basis
for comparison, while Sections~\ref{results_sect}, ~\ref{discussion_sect},
and ~\ref{conclusions_sect} give results, discussion, and conclusion.

\section{The Coulomb Free Energy \lab{coulomb_sect}}

The general expression adopted for the Coulomb free energy employs the
following definitions for the Coulomb sums $\Sigma^{\mboxs{C}}_0$ and
$\Sigma^{\mboxs{C}}_2$, the weak Coulomb interaction parameter $\Lambda$,
and the strong Coulomb interaction parameter $\Gamma$: \be
\Sigma^{\mboxs{C}}_0 \equiv \sum_\Pi n_\Pi, \lab{sig0_eq1} \ee \be
\Sigma^{\mboxs{C}}_2 \equiv \sum_\Pi n_\Pi {Z_\Pi}^2, \lab{sig2_eq1} \ee \be
\Lambda \equiv \frac{2 e^3 \pi^{1/2}(\Sigma^{\mboxs{C}}_2 + n_e
\theta_e)^{3/2}} {(4 \pi \epsilon_0 kT)^{3/2}(\Sigma^{\mboxs{C}}_0 + n_e
\theta_e)}, \lab{lambda_eq}\ee and \be \Gamma \equiv (\Lambda/3^{1/2})^{2/3}
= \frac{e^2}{4 \pi \epsilon_0 kT} \frac{(4 \pi/3)^{1/3}\Sigma^{\mboxs{C}}_2
+ n_e \theta_e}{(\Sigma^{\mboxs{C}}_0 + n_e \theta_e)^{2/3}}.
\lab{gamma_eq}\ee The number density of the positive ion identified by index
$\Pi$ is given by $n_\Pi$, the charge on that positive ion is given by
$Z_\Pi$, and \be \theta_e \equiv \partial \ln n_e(\eta,T)/\partial \eta,
\lab{theta_eq}\ee where $n_e$ is the number density of free electrons and
$\eta$ is a degeneracy parameter.  I describe the approximation used for
$n_e(\eta,T)$ in Paper~I, and $\theta_e$ is calculated from the analytical
derivative of that approximation. (Note all units in this paper are SI
unless specified otherwise.  Furthermore, throughout this paper the symbol
$\ln$ will stand for the natural logarithm and the symbol $\log$ will stand
for the base 10 logarithm.) Using the above definitions, the adopted
expression for the Coulomb free energy reduces to \be F_{\mboxs{C}} = -kTV
(\Sigma^{\mboxs{C}}_0 + n_e \theta_e) g_{\mboxs{C}}(\Gamma), \lab{fc_eq} \ee
where $g_{\mboxs{C}}(\Gamma)$ is the so-called Coulomb function.

These expressions for $\Lambda$, $\Gamma$, and $F_{\mboxs{C}}$ are the same
as those used by Pols et al. (1995, PTEH), except we have
replaced $\Sigma^C_0$ in their formulation by $\Sigma^C_0 + n_e \theta_e$.
In the limit of small degeneracy (which correlates with small $\Gamma$),
$\theta_e$ approaches unity, and our formulation reduces to the DeWitt (1969)
formulation.  In the limit of large degeneracy (which correlates with large
$\Gamma$), $\theta_e$ approaches zero and our formulation reduces to the
PTEH formulation.  The recommended Coulomb function,
$g_{\mboxs{C}}(\Gamma)$, is discussed in the following sub-sections.

\subsection{Constraints on the recommended Coulomb function and associated free
energy \lab{constraint_sect}}

For small $\Gamma$ the recommended Coulomb function should approach the
Debye-H\"uckel (DH) result, \be g_{\mboxs{DH}}(\Gamma) \equiv
\Gamma^{3/2}/3^{1/2}. \lab{gdh_eq} \ee This small-$\Gamma$ constraint on the
Coulomb function insures the corresponding recommended Coulomb free-energy
will approach \be F_{\mboxs{DH}} = -kTV (\Sigma^{\mboxs{C}}_0 + n_e
\theta_e) \Lambda/3, \lab{fdh_eq}\ee which for small degeneracy gives the
same result that can be derived from from first principles using an activity
expansion of the grand canonical partition function (Rogers 1994).

For large $\Gamma$ the recommended Coulomb function should approach the
one-component-plasma (OCP) result which is well represented by \be
g_{\mboxs{OCP}}(\Gamma) = a\Gamma + b\Gamma^s + c \ln \Gamma + d,
\lab{gocp_eq}\ee where $a=0.899126$, $b=-0.60712/s$, $s=0.321308$,
$c=0.27998$, and $d = 0.436484 - a - b$ (see DeWitt, Slattery, \& Chabrier
1996, DSC). This large-$\Gamma$ constraint on the Coulomb function
insures the corresponding Coulomb free-energy approximates liquid,
multi-component-plasma (LMCP) results which are well represented by the
linear mixing rule, \be F_{\mboxs{LMCP}} = -kTV \sum_\Pi n_\Pi
g_{\mboxs{OCP}}(Z_\Pi^{5/3}\Gamma_0), \lab{linear_mixing_rule_eq}\ee where
\be \Gamma_0 = \frac{e^2}{4 \pi \epsilon_0 kT}\left (\frac{4 \pi}{3}
\Sigma^{\mboxs{C}}_1 \right ) ^{1/3}\ee and \be \Sigma^{\mboxs{C}}_1 \equiv
\sum_\Pi n_\Pi Z_\Pi. \lab{sig1_eq1} \ee For large degeneracy, $\theta_e
\approx 0$; and for large $\Gamma$, $g_{\mboxs{OCP}} \approx a \Gamma$.  Under
these conditions, the relative difference between the Coulomb free energy
corresponding to $g_{\mboxs{OCP}}$ and the LMCP result is \be
\frac{\Sigma^{\mboxs{C}}_0 \Gamma}{\sum_\Pi n_\Pi Z_\Pi^{5/3}\Gamma_0} - 1 =
\frac{(\Sigma^{\mboxs{C}}_0)^{1/3}\Sigma^{\mboxs{C}}_2}
{(\Sigma^{\mboxs{C}}_1)^{1/3} \Sigma^{\mboxs{C}}_{5/3}} - 1,
\lab{large_gamma_error_eq}\ee where \be \Sigma^{\mboxs{C}}_{5/3} \equiv
\sum_\Pi n_\Pi {Z_\Pi}^{5/3}. \lab{sig5/3_eq1} \ee Note both
equation~(\ref{gocp_eq}) and equation~(\ref{linear_mixing_rule_eq})are not
exact representations of the definitive DSC Monte Carlo results, but I have
ignored these representation errors in this development since, in general,
they are small compared to the relative error given by
equation~(\ref{large_gamma_error_eq}).

\subsection{The Recommended Coulomb Function\lab{gc_sect}}

In the FreeEOS implementation the recommended Coulomb function is
represented by \be \log g_{\mboxs{C}} \equiv Y(X), \ee where $X \equiv \log
\Gamma$ and $Y(X)$ is characterized by the two fitting parameters
$X_{\mboxs{DH}}$ and $X_{\mboxs{OCP}}$.  For all $X$ values less than or
equal to $X_{\mboxs{DH}}$, $Y$ is set equal to the DH result, $\log
g_{\mboxs{DH}}$, which more than satisfies the small-$\Gamma$ constraint on
$g_{\mboxs{C}}$ discussed in Section~\ref{constraint_sect}.  For all $X$
values greater than or equal to $X_{\mboxs{OCP}} > X_{\mboxs{DH}}$, $Y$ is
set equal to the modified OCP result, $\log (g_{\mboxs{OCP}} + \delta c \ln
\Gamma + \delta d)$, which satisfies the large-$\Gamma$ constraint on
$g_{\mboxs{C}}$ discussed in Section~\ref{constraint_sect}.
$X_{\mboxs{DH}}=-0.4$ and $X_{\mboxs{OCP}}=0.$ give the best fit of FreeEOS
results to OPAL results (see Sect.~\ref{results_sect}) and these values are
adopted for the recommended Coulomb function.

The parameters $\delta c$ and $\delta d$ are internally adjusted in the
FreeEOS implementation to give a cubic polynomial in the transition region
between $X_{\mboxs{DH}}$ and $X_{\mboxs{OCP}}$ that is continuous in $Y(X)$,
$Y^{\prime}(X)$, and $Y^{\prime\prime}(X)$ with DH results at
$X=X_{\mboxs{DH}}$ and with modified OCP results at $X=X_{\mboxs{OCP}}$. The
basic equations that must be satisfied by that cubic are given in Sect.3.3
of Press et al. (1986).  The cubic is determined by the 4 requirements of
$Y(X)$ and $Y^{\prime\prime}(X)$ continuity at $X=X_{\mboxs{DH}}$ and
$X=X_{\mboxs{OCP}}$. However, for the unmodified OCP result (i.e., $\delta
c=0$ and $\delta d=0$), the first derivatives are non-continuous.  To make
them continuous, a Newton-Raphson iteration procedure is used to solve the
two $Y^{\prime}(X)$ continuity constraints in terms of the two unknowns
$\delta c$ and $\delta d$.  This procedure is only done once for a
particular set of $X_{\mboxs{DH}}$ and $X_{\mboxs{OCP}}$ values, and the
resulting $\delta c$ and $\delta d$ values tend to be small so that modified
OCP results from the recommended Coulomb function do not strongly differ
from unmodified OCP results for large $\Gamma$ values (see
Sect.~\ref{results_sect}).

\subsection{Other Models for the Coulomb Free Energy}

The recommended treatment of the Coulomb free energy (i.e.,
eq.~[\ref{fc_eq}] with the Coulomb function, $g_{\mboxs{C}}(\Gamma)$,
constructed as in Sect.~\ref{gc_sect}) has been designed to join the
Debye-H\"uckel approximation smoothly with a good approximation to LMCP
results using the Coulomb fitting parameters $X_{\mboxs{DH}} = -0.4$ and
$X_{\mboxs{OCP}} = 0.$ This treatment is part of the EOS1 free-energy model
option suite (described in Paper~I) that gives a good fit to OPAL solar
results (see Paper~II). However, I have implemented additional Coulomb
free-energy model options for FreeEOS. These options provide variations on
the EOS1 option suite that reduce the required computer time to perform the
EOS calculations.  Also these options (in combination with other FreeEOS
options) provide almost exact emulation of the PTEH; Swenson, Irwin, \&
Rogers, (unpublished, SIREFF); or Mihalas, D\"{a}ppen, \& Hummer (1988, MDH)
equations of state.

\subsubsection{Approximations for the Coulomb Sums \lab{csum_sect}}

The form of the Coulomb free energy (eq.~[\ref{fc_eq}]) implies the
associated correction to the equilibrium constants given by equation~(35) of
Paper~II depends on the auxiliary variables $\Sigma^{\mboxs{C}}_0$ and
$\Sigma^{\mboxs{C}}_2$ defined by equations~(\ref{sig0_eq1}) and
(\ref{sig2_eq1}).  Normally, these auxiliary variables are determined to
high numerical precision as part of the EOS solution by Newton-Raphson
iteration as discussed in Paper~II, but the FreeEOS implementation also
offers two options for approximating these auxiliary variables.  These
approximations save computer time by reducing the number of auxiliary
variables that need to be determined by iteration.

The ``PTEH'' approximations for the Coulomb sums are taken from the PTEH
reference, and are given as follows: \be \Sigma^{\mboxs{C}}_0 \approx n_e
\Sigma^+_0/\Sigma^+_1 \lab{sig0_eq3} \ee and \be \Sigma^{\mboxs{C}}_2
\approx n_e \Sigma^+_2/\Sigma^+_1 ;\lab{sig2_eq3} \ee where for $k = 0,1,2$
\be \Sigma^+_k = \sum_j \epsilon_j Z_j^k. \ee The $j$ index runs over all
elements, $Z_j$ is the nuclear charge (atomic number), and \be \epsilon_j =
X_j/A_j. \ee $X_j$ is the abundance by weight, $A_j$ is the atomic weight,
and $\sum_j X_j \equiv 1$.  The approximations given by
equations~(\ref{sig0_eq3}) and (\ref{sig2_eq3}) are exact if all elements
are un-ionized or completely ionized.  Furthermore, these approximations
require evaluation of no auxiliary variables.

The ``SIREFF'' approximations for the Coulomb sums are taken from the SIREFF
EOS implementation and are given as follows: \be \Sigma^{\mboxs{C}}_0
\approx n(H^+) + n(He^+) + n(He^{++}) +
\frac{n(\mboxs{H}^+)}{\epsilon({\mboxs{H}})} \sum_{m}
\epsilon_{m}\lab{sig0_eq2} \ee and \be \Sigma^{\mboxs{C}}_2 \approx n(H^+) +
n(He^+) + 4 n(He^{++}) + \frac{n(\mboxs{H}^+)}{\epsilon({\mboxs{H}})}
\sum_{m} \epsilon_{m} + \frac{n(\mboxs{He}^{++})} {\epsilon({\mboxs{He}})}
\sum_{m} \epsilon_{m}(Z_{m}^2 - 1); \lab{sig2_eq2} \ee where the $m$ index
runs over just the metals. Equations~(\ref{sig0_eq2}) and (\ref{sig2_eq2})
are exact in the limits where all elements are un-ionized, singly ionized or
completely ionized. These approximations normally require knowledge of no
additional auxiliary variables when using the SIREFF free-energy model.  For
that model, $n(\mboxs{H}_2^+)$ is ignored (which is why that species is not
included in the above equations), and $n(\mboxs{H}^+)$, $n(\mboxs{He}^+)$,
and $n(\mboxs{He}^{++})$ are already required auxiliary variables for the
SIREFF form of pressure ionization.

\subsubsection{Additional Options for the Coulomb Treatment
\lab{coulomb_compare_sect}}

Other Coulomb options available for FreeEOS users are as follows.
\begin{description}
\item[(1)] Ignore the Coulomb effect.
\item[(2)] Replace $X_{\mboxs{DH}} = -0.4$ and $X_{\mboxs{OCP}} = 0.$ by
$X_{\mboxs{DH}} = -1.$ and $X_{\mboxs{OCP}} = 0$.
\item[(3)] Replace $\Sigma^{\mboxs{C}}_0 + n_e \theta_e$ wherever it
appear in the recommended formulation by $\Sigma^{\mboxs{C}}_0$.
\item[(4)] Replace the recommended Coulomb function by the PTEH version
(see their eq.~[25]), \be g_{\mboxs{PTEH}}(\Gamma) = \frac{a_1
\Gamma}{\{[\Gamma/(\Gamma+a_3)]^{1/a_2} +
(a_1\sqrt{3/\Gamma})^{1/a_2}\}^{a_2}}, \lab{gpteh_eq}\ee where
$a_1=0.89752$, $a_2=0.768$, and $a_3=0.208$.
\item[(5)] Replace the recommended $\theta_e$ (see discussion following
eq.~[\ref{theta_eq}]) by the PTEH version (which is based on their
eq.~[29]).
\item[(6)] Replace the recommended Coulomb function by the Debye-H\"uckel
version $g_{\mboxs{DH}}$ (eq.~[\ref{gdh_eq}]).
\item[(7)] Replace the recommended Coulomb function by a corrected
Debye-H\"uckel version $\tau(x) g_{\mboxs{DH}}$; where \be \tau(x) =
3x^{-3}[\ln(1+x)-x+x^2/2] \lab{tau_eq}\ee (eq.~[11] from MDH), \be x \equiv
\left [ \frac{2 \pi^{1/2} e^3 n_e}{(4 \pi \epsilon_0 kT)^{3/2}} \right ]
\frac{F_{1/2}(\eta,0)}{F_{3/2}(\eta,0)} \frac{(\Sigma^{\mboxs{C}}_2 +
\theta_e n_e)^{1/2}}{\Sigma^{\mboxs{C}}_0}, \lab{x_eq}\ee and
$F_{k}(\eta,0)$ is the non-relativistic limit of the Fermi-Dirac integral of
order $k$ (see Paper~I).
\end{description}

Options~(1) and (2) above are useful for studies (e.g., Cassisi, Salaris, \&
Irwin 2003) of the sensitivity of stellar-interior results to the overall
Coulomb effect and variations in the approximation for the Coulomb effect at
intermediate $\Gamma$; the combination of the PTEH approximations for the
Coulomb sums and options~(3) through (5) above constitute the PTEH Coulomb
model; the combination of the SIREFF approximations for the Coulomb sums and
option~(6) above constitute the SIREFF Coulomb model; and precise iterative
calculation of the Coulomb sums and option~(7) above constitute the MDH
Coulomb model.

\section{Results \lab{results_sect}}

Figure~\ref{coulomb_adjust_fig} gives the recommended $\log g_{\mboxs{C}}$;
its first two derivatives; and residuals of $\log g_{\mboxs{DH}}$, $\log
g_{\mboxs{OCP}}$, and $\log g_{\mboxs{PTEH}}$ relative to $\log
g_{\mboxs{C}}$.  By design, the recommended $g_{\mboxs{C}}$ and its first
and second derivative are continuous and agree with $g_{\mboxs{DH}}$ for
$\log \Gamma \le X_{\mboxs{DH}} = -0.4$.  Because of the continuity
constraint discussed in Sect.~\ref{gc_sect}, $\delta c$ and $\delta d$ are
non-zero so that $\log g_{\mboxs{C}}$ differs from $\log g_{\mboxs{OCP}}$ by
0.095 at $X=0.$, but these differences decrease with increasing $X$.  The
differences between $g_{\mboxs{C}}$ and $g_{\mboxs{PTEH}}$ are discussed in
the next section.

To give the following results and discussion some context,
Figure~\ref{context_fig} compares EOS quantities and the loci of several
stellar interior models on the density-temperature plane for solar
metallicity. The high-density, low-temperature calculation limit indicated
in the figure is defined by \be \log \rho_{\mboxs{lim}} = \log \rho_5 +
(3/2) \log(T/10^5)\ee for $\log T < 6$ and continued by the $\log T = 6$
isotherm. The limit parameter of $\log \rho_5 = 3.3$ is used for this figure
to avoid regions where FreeEOS calculations currently do not converge. This
calculation limit roughly corresponds to a 0.1-M$_\sun$ model and can also
be viewed as the approximate limit of validity for FreeEOS calculations (see
further discussion in Paper~II).  This same $\log \rho_5$ limit parameter is
also used for all other $\rho,T$ figures in this paper except for
Figure~\ref{dh_fig}.

Figure~\ref{nocoulomb_fig} shows the overall change in pressure due to the
Coulomb effect.  In general, the Coulomb effect increases with increasing
density and decreasing temperature, but this overall trend is modified by
other factors which affect the EOS. The Coulomb effect is reduced at high
densities because the pressure from degenerate electrons grows faster than
the Coulomb pressure, at high temperatures because the ideal pressure
increases with temperature while the Coulomb pressure decreases, combined
high temperatures and low densities because the radiation pressure dominates
the pressure, at low densities because the mean distance between particles
is low (and therefore interactions are reduced), and at low temperatures
because of low ionization (fewer charged particles to interact). The net
result of all these different effects is as illustrated; for fixed abundance
the Coulomb effect is largest at combined intermediate densities and
intermediate temperatures corresponding to the envelopes of extreme
lower-main-sequence stars. In general, for fixed abundance and any given
stellar model the Coulomb effect first increases inward as ionization is
increased, reaches a maximum in the envelope, and then decreases inward as
the temperature is increased still further.  Note in contrast to the fixed
abundance case, the Coulomb effect may increase inward near the core of
models in advanced stages of stellar evolution because the effect goes as
the cube of the charge on the ions (see eqs.~[\ref{lambda_eq}] and
[\ref{fdh_eq}]) which monotonically increases as a function of time due to
nuclear processing.

Figure~\ref{dh_fig} shows the pressure effect of changing the model for the
Coulomb function from the recommended $g_{\mboxs{C}}$ to $g_{\mboxs{DH}}$.
The two FreeEOS calculations being compared are identical by definition for
$\log \Gamma \le X_{\mboxs{DH}} = -0.4$ (see Fig.~\ref{coulomb_adjust_fig}),
but as the $\Gamma$ value is increased the two calculations rapidly diverge
to the point where the FreeEOS calculation that uses $g_{\mboxs{DH}}$
produces unphysical results (e.g., negative overall pressures) for physical
conditions that are typical for envelopes of extreme LMS stars. This
illustrates why a non-DH Coulomb approach is required for such conditions.

Figures~\ref{pteh_fig} through \ref{0.1_model_pteh_fig} show the effect of
changing the model for the Coulomb function from the recommended
$g_{\mboxs{C}}$ to $g_{\mboxs{PTEH}}$.  Figure~\ref{pteh_fig} shows the
overall change in pressure, while figures~\ref{1.0_model_pteh_fig} and
\ref{0.1_model_pteh_fig} show more detailed results for the locus of points
in a solar-mass model and a 0.1-$M_\sun$ model.  The model results will be
discussed in detail in the next section.

\section{Discussion \lab{discussion_sect}}

The Coulomb functions $g_{\mboxs{C}}$ and $g_{\mboxs{PTEH}}$ are continuous
in at least the function and its first and second derivatives.  This is a
necessary condition for second-order thermodynamic quantities like
$\Gamma_1$ (which depend on a second-order partial derivative of the
free-energy) to be continuous.  Furthermore, $g_{\mboxs{C}}$ and
$g_{\mboxs{PTEH}}$ satisfy the small- and large-$\Gamma$ constraints
presented in Section~\ref{constraint_sect}. Thus, the EOS results
corresponding to these Coulomb functions make a smooth transition from DH
results at small $\Gamma$ to an approximation to LMCP results at large
$\Gamma$.  However, Figure~\ref{coulomb_adjust_fig} shows that
$g_{\mboxs{PTEH}}$ has significant differences with $g_{\mboxs{C}}$. For
$\log \Gamma \le X_{\mboxs{DH}} = -0.4$, these differences are identical to the
differences of $g_{\mboxs{PTEH}}$ relative to $g_{\mboxs{DH}}$ while for
$\log \Gamma \ge X_{\mboxs{OCP}} = 0.$, these differences are essentially the
negative of the differences between $g_{\mboxs{C}}$ and $g_{\mboxs{OCP}}$.

Figure~\ref{1.0_model_pteh_fig} presents FreeEOS results as a function of
the model depth variable $T(r)$ in a solar-mass model. The first panel of
the figure shows how $\Gamma_1$ behaves with depth to indicate important EOS
transitions such as the combined hydrogen and first helium ionization zone
(the primary minimum in $\Gamma_1$) and the the second helium ionization
zone (the secondary minimum in $\Gamma_1$).  The remaining panels of the
figure show how the differences between $g_{\mboxs{PTEH}}$ and
$g_{\mboxs{C}}$ (and essentially the differences between $g_{\mboxs{PTEH}}$
and $g_{\mboxs{DH}}$ since the maximum solar $\log \Gamma$ value is only
slightly larger than $X_{\mboxs{DH}} = -0.4$) propagate to various
thermodynamic quantities for the locus of points in a solar-mass model.

To help predict the effect of the illustrated residuals on derived solar
properties it is important to remember that in the adiabatic part of the
solar convection zone, $P$ becomes an implicit function of $\rho$ along an
adiabat.  Once the adiabatic exponent, $\Gamma_1 = \partial \ln P(\rho,
s)/\partial \ln \rho$ is known along the adiabat (defined by constant
entropy per unit mass, $s$), the solution for the density-depth and
pressure-depth relations, $\rho(r)$ and $P(r)$ follows from the boundary
conditions on the adiabatic region and solution of Poisson's equation (see
discussion in Cox and Giuli 1968, Sect.~19.1).  The temperature-depth
relation, $T(r)$ is then determined from the EOS relation $T(\rho(r), P(r))$
in the adiabatic region. The residuals in $\Gamma_1$ tend to be oscillatory
as illustrated in the figure because $\Gamma_1$ is a second-order
thermodynamic quantity.  These oscillatory residuals should have little
effect on $\rho(r)$ and $P(r)$ outside the adiabatic region. In contrast,
the EOS relation $P(\rho(r),T(r))$ is a first-order thermodynamic quantity
(it depends on a first partial derivative of the free energy) so $\ln
P(\rho(r),T(r))$ residuals tend to be systematic as illustrated in the
figure.  The residuals in the $\ln T(\rho(r), P(r))$ relation should
similarly be systematic (although opposite in sign). So long as these
systematic residuals as illustrated in this figure do not propagate to the
surface or to the deep part of the model where the luminosity is generated,
and so long as the residuals in $\ln \rho(r)$ and $P(r)$ do not propagate
outside the adiabatic zone, then the radius and luminosity of the model
should largely be unaffected by EOS changes or uncertainties in the
adiabatic part of the model.

In contrast to the case of the predicted radius and luminosity of the sun,
the predicted acoustic periods of the sun should be significantly affected
(at least compared to the high level of accuracy with which the solar
acoustic periods have been observationally determined) by the difference
between $g_{\mboxs{PTEH}}$ and $g_{\mboxs{DH}}$.  The predicted acoustic
periods are largely determined by the integral of the inverse sound speed
over depth, and the change in adopted Coulomb function affects the sound
speed at the $10^{-2}$ level right where the integrand is fairly near its
maximum value. Furthermore, comparison of this plot with Figures~6 and 7 of
Paper~II, shows use of $g_{\mboxs{PTEH}}$ completely destroys the good
agreement we have with the OPAL EOS for the solar case. In fact, in all the
attempts to adjust the EOS1 option suite to fit the OPAL results, I have
never been able to improve on simply adopting the Debye-H\"uckel Coulomb
function, $g_{\mboxs{DH}}$, for the solar case.  It is for this reason that
I have designated $g_{\mboxs{C}}$ (which is essentially identical to
$g_{\mboxs{DH}}$ in the solar case by design) as the recommended Coulomb
function and use that function rather than $g_{\mboxs{PTEH}}$ in the EOS1
option suite.

Figure~\ref{0.1_model_pteh_fig} shows the same quantities as
Figure~\ref{1.0_model_pteh_fig} except that all calculations have been done
for a 0.1-$M_\sun$ model.  The envelope conditions of this model are near
the limit of validity of the current EOS1 option suite of FreeEOS because of
deficiencies in the adopted pressure ionization scheme.  In fact, the
current pressure-ionization model for the metals is much less detailed than
it is for hydrogen and helium, and this oversimplification causes
near-discontinuous metal ionization and the accompanying sharp changes that
occur for $\Gamma_1$ and other quantities (for example, near $\log T(r) =
5.7$).

The Coulomb effect is large ($\log \Gamma > 0$) through much of the
0.1-$M_\sun$ model with the maximum $\log \Gamma$ value approaching 0.5 in
the envelope.  Thus, through much of the model $g_{\mboxs{PTEH}}$ will give
a better approximation to $g_{\mboxs{OCP}}$ than $g_{\mboxs{C}}$, and a
FreeEOS calculation with $g_{\mboxs{PTEH}}$ will give a better approximation
to LMCP results than the $g_{\mboxs{C}}$ calculation.  Nevertheless, except
for the outermost layers the model is completely convective, and because of
the high densities involved almost all but the outermost layers are
adiabatic as well. Under these conditions (as discussed for the solar-mass
model), the largely oscillatory nature of the $\Gamma_1$ residuals and lack
of systematic $\ln P(\rho(r),T(r))$ residuals at the surface and at depth
indicate that the predicted radius and luminosity of the 0.1-$M_\sun$ model
should largely be unaffected by using $g_{\mboxs{C}}$ rather than
$g_{\mboxs{PTEH}}$.  An even stronger argument can be made for the relative
error of the current Coulomb formulation compared to the linear-mixing rule
formulation for large $\Gamma$.  Indeed, for normal helium abundances the
maximum $n(\mboxs{He}^{++})/n(\mboxs{H}^{+})$ value is approximately 0.10,
and the maximum error estimated from equation~\ref{large_gamma_error_eq} is
only 3~\% which is an order of magnitude smaller than the relative
differences between $g_{\mboxs{PTEH}}$ and $g_{\mboxs{C}}$ shown in
Figure~\ref{coulomb_adjust_fig}.

\section{Conclusions \lab{conclusions_sect}}

This paper has described how the Coulomb correction is implemented in the
FreeEOS software library. The recommended set of Coulomb options that are
used in the EOS1 option suite smoothly join Debye-H\"uckel results at weak
Coulomb coupling to a good approximation to liquid, multi-component-plasma
results for strong Coulomb coupling.  The match with the Debye-H\"uckel
results is maintained to a sufficiently large value of the Coulomb coupling
constant ($\log \Gamma = -0.4$) to secure an excellent fit to OPAL results
for the solar case (see Figs.~6 and 7 of Paper~II). This Coulomb treatment
should be suitable for most stellar-interior calculations with the notable
exception of the cool white dwarfs where the phase transition to a
crystalline state is encountered for large Coulomb coupling.

The FreeEOS software library is licensed under the GPL and is freely
downloadable from \url{http://freeeos.sourceforge.net/}. Older versions of
FreeEOS used a spline fit to approximate the Coulomb implementation
introduced at version 1.4.0 of FreeEOS and described in this paper.

\acknowledgments I thank Forrest Rogers for many useful discussions; Santi
Cassisi for providing some representative model calculations and for his
friendly encouragement of my FreeEOS work; Ben Dorman, Fritz Swenson, and
Don VandenBerg for helping to arouse my original interest in the EOS problem
for stellar interiors; and Richard Stallman, Linus Torvalds, and many other
programmers for the GNU/Linux computer operating system and accompanying
tools that have made it practical to develop the FreeEOS code on personal
computers. The figures of this paper have been generated with the
yPlot~(\url{http://yplot.sourceforge.net}) and
PLplot~(\url{http://www.plplot.org}) scientific plotting packages.

\clearpage
\begin {references}
\reference{Cass2003}
Cassisi, S., Salaris, M., \& Irwin, A.W. 2003, \apj, 588, 862
\reference{Cass2005}
Cassisi, S. 2005, private communication
\reference{Cox1968}
Cox, J. P., \& Giuli, R. T. 1968, Principles of Stellar Structure (New York:
Gordon and Breach)
\reference{DeWi1969}
DeWitt, H. E. 1969, in Low Luminosity Stars, ed. S. Kumar (New York:
Gordon and Breach), p. 211
\reference{DeWi1996}
DeWitt, H., Slattery, W., \& Chabrier, G. 1996, Physica B228, 21 (DSC)
\reference{Irwi2004a}
Irwin, A.W. 2004a, \url{http://freeeos.sourceforge.net/eff\_fit.pdf} (Paper~I)
\reference{Irwi2004b}
Irwin, A.W. 2004b, \url{http://freeeos.sourceforge.net/solution.pdf} (Paper~II)
\reference{Miha1988}
Mihalas, D., D\"{a}ppen, W., \& Hummer, W. G. 1988, \apj, 331, 815 (MDH)
\reference{Piet2004}
Pietrinferni, A., Cassisi, S., Salaris, M., \& Castelli, F. 2004, \apj,
612, 168
\reference{Pols1995}
Pols, O. R., Tout, C. A., Eggleton, P. P., \& Han, Z. 1995,
\mnras, 274, 964 (PTEH)
\reference{Pres1986}
Press, W.H., Flannery, B.P., Teukolsky, S.A., \& Vetterling, W. T. 1986,
Numerical Recipes (Cambridge, Cambridge University Press)
\reference{Roge1994}
Rogers, F. J. 1994, in The Equation of State in Astrophysics IAU Coll. 147,
ed. G. Chabrier, E. Schatzman (Cambridge: Cambridge University Press), p. 16
\reference{Roge2002}
Rogers, F.J. \& Nayfonov, A. 2002, \apj, 576, 1064 (OPAL)

\end{references}

\clearpage
\begin{figure}
\plotone{coulomb_adjust.\FigureSuffix}
\figcaption{The recommended $Y(X) = \log g_{\mboxs{C}}$ and its first two
derivatives as a function of $X = \log \Gamma$; and a comparison of $\log
g_{\mboxs{DH}}$, $\log g_{\mboxs{OCP}}$, and $\log g_{\mboxs{PTEH}}$ with
$Y$ as a function of $X$. $Y(X)$ smoothly joins DH results with modified OCP
results using a cubic polynomial (see text). \lab{coulomb_adjust_fig}}
\end{figure}

\clearpage
\begin{figure}
\plotone{context.\FigureSuffix}
\figcaption{A solar-metallicity comparison of EOS quantities and loci of
model stellar interiors. The thin short-dashed line indicates where the
radiative and gas pressures are equal and is the boundary between the
radiation-dominated and matter-dominated EOS.  The thin solid lines indicate
the middle of important ionization and dissociation zones. The ``He$^+$'',
``He'', ``H'', and ``H$_2$'' labels respectively correspond to where half
the helium is singly ionized or neutral and where half the hydrogen is in
neutral monatomic form or in the diatomic molecular form. The medium
thickness solid line indicates the current high-density, low-temperature
calculation limit for FreeEOS (see text). The thick solid lines indicate
$\rho,T$ loci of model stellar interiors provided by Cassisi (2005).  These
models were calculated using the EOS1 option suite of FreeEOS and the
stellar-evolution code that has been described in Pietrinferni, et al.
(2004). The labels of ``0.1'', ``0.3'', ``1.0'', ``RGT'', and ``CG''
respectively indicate main-sequence models of 0.1, 0.3, and 1.0 $M_{\sun}$
and models of 1.0 $M_{\sun}$ evolved to the tip of the red-giant branch and
to the initial clump-giant phase (zero-age horizontal branch of solar
metallicity). \lab{context_fig}}
\end{figure}

\clearpage
\begin{figure}
\plotone{nocoulomb.\FigureSuffix}
\figcaption{The log of the relative change in pressure caused by changing
from the recommended free-energy model to one where the Coulomb free-energy
component is dropped. Solar abundances were used for both EOS calculations
employed in the comparison. The high-density, low-temperature calculation
limit is the same as in Figure~\ref{context_fig}. \lab{nocoulomb_fig}}
\end{figure}

\clearpage
\begin{figure}
\plotone{dh.\FigureSuffix}
\figcaption{The log of the relative change in pressure caused by changing
from the recommended Coulomb free-energy model to one where $g_{\mboxs{C}}$
is replaced by $g_{\mboxs{DH}}$ (eq.~[\ref{gdh_eq}]). Solar abundances were
used for both EOS calculations employed in the comparison. To avoid negative
pressures for the DH case, the high-density, low-temperature calculation
limit was shifted to lower densities by 1.05 dex (i.e., $\log \rho_5$ was
reduced from 3.3 to 2.25). \lab{dh_fig}}
\end{figure}

%\clearpage
%\begin{figure}
%\plotone{dhtau.\FigureSuffix}
%\end{figure}

\clearpage
\begin{figure}
\plotone{pteh.\FigureSuffix}
\figcaption{The log of the relative change in pressure caused by changing
from the recommended Coulomb free energy model to one where $g_{\mboxs{C}}$
is replaced by $g_{\mboxs{PTEH}}$ (eq.~[\ref{gpteh_eq}]).  Solar abundances
were used for both EOS calculations employed in the comparison. The
high-density, low-temperature calculation limit is the same as in
Figure~\ref{context_fig}. \lab{pteh_fig}}
\end{figure}

\clearpage
\begin{figure}
\plotone{{1.0_pteh}.\FigureSuffix}
\figcaption{The adiabatic pressure gradient, $\Gamma_1$, for the recommended
Coulomb free-energy model as well as the the change in $\Gamma_1$, relative
change in $P$, and the log of the relative change in the square of the sound
speed, ${v_s}^2 = \Gamma_1 P/\rho$, as a function of the temperature of each
depth point for a solar-mass model.  The plotted changes are derived from
two independent FreeEOS calculations; one with the recommended Coulomb
treatment and one where $g_{\mboxs{C}}$ has been replaced by
$g_{\mboxs{PTEH}}$ (eq.~[\ref{gpteh_eq}]). \lab{1.0_model_pteh_fig}}
\end{figure}

\clearpage
\begin{figure}
\plotone{{0.1_pteh}.\FigureSuffix}
\clearpage
\figcaption{The same as figure~\ref{1.0_model_pteh_fig} except all
calculations were done for a 0.1-M$_\sun$ model. \lab{0.1_model_pteh_fig}}
\end{figure}

\end{document}
